\chapter{Tensor Calculus, Christoffel Symbols, and Geodesics}

In this lecture we move from tensor algebra to tensor calculus. The issue is no longer merely how tensors add or transform, but how to express physics so that the truth of an equation does not depend on coordinates. That question forces us to confront an obstruction: ordinary derivatives of tensor components do not in general transform tensorially. The lecture builds the remedy slowly. We begin with scalars, where the ordinary derivative behaves well, pass through the failure for vectors, extract the extra term that measures the flopping around of the coordinates themselves, and then turn that nuisance into the Christoffel symbol and the covariant derivative. From there the discussion broadens to the metric, then to differentiation along curves, and finally to geodesics, which are the straightest possible curves in a curved space.

\section{Why Tensor Calculus Enters General Relativity}

We begin with a brief reminder from tensor algebra. Compatible tensors may be added and subtracted: one may add covariant components to covariant components, or contravariant components to contravariant components, but not mismatched objects. The point of the reminder is not algebra for its own sake. It is to clear the ground for the real question: why do we care about tensors at all?

The answer is that the statements of physics ought to survive a change of coordinates. The numbers on the left-hand side and right-hand side of an equation may change when we pass from one chart to another, but the fact that two sides are equal should not. Equivalently, if we can write a law in the form that some object is zero, then we want zero in one coordinate system to imply zero in every coordinate system.

This is exactly the virtue of tensors. If a tensor vanishes in one frame, it vanishes in all frames. To see the idea in the simplest cases, recall the transformation law for a covariant vector:
\begin{equation}
V_n(y)=\frac{\partial x^r}{\partial y^n}V_r(x).
\end{equation}
If every component \(V_r(x)\) is zero, then every component \(V_n(y)\) is zero as well. Likewise for a mixed tensor,
\begin{equation}
T^m{}_n(y)=\frac{\partial y^m}{\partial x^r}\frac{\partial x^s}{\partial y^n}T^r{}_s(x),
\end{equation}
and again vanishing in one coordinate system implies vanishing in every other.

So the aim of tensor calculus is clear. We want a notion of differentiation that preserves this tensorial character. General relativity will need such an operation because physics is not only about what tensors are, but how they vary from point to point.

\section{Scalars First, Then the Trouble with Vectors}

The lecture wisely begins with scalars. A scalar does not transform at all; it is just the same number attached to the same geometric point. If \(\phi\) is a scalar field, then the chain rule gives
\begin{equation}
\frac{\partial \phi}{\partial y^n}
=
\frac{\partial x^m}{\partial y^n}
\frac{\partial \phi}{\partial x^m}.
\end{equation}
Thus the derivative of a scalar transforms as a covariant vector. This is the easy case.

There is also a physical way to say the same thing. If the temperature in a room is constant everywhere, then it remains constant no matter what coordinates we choose. Polar coordinates, Cartesian coordinates, or any more exotic system will all agree that the temperature has no spatial variation. So when the derivative of a scalar vanishes, it makes sense that it vanishes in every coordinate system.

Now we come to vectors, and here the difficulty appears. Suppose in ordinary flat space we have a vector field that points in the same direction everywhere and has the same length everywhere. In Cartesian coordinates its components are constant, so their ordinary derivatives vanish. But that conclusion depends on the coordinates being rigid.

If instead we lay curvilinear coordinates on the same flat blackboard, the axes themselves rotate and bend from point to point. The vector may remain geometrically the same while its components change, because components depend not only on the vector but also on the local orientation of the coordinate axes. Thus a geometrically constant vector field need not have constant components in curvilinear coordinates.

That is the first warning. Ordinary derivatives of vector components do not behave like the derivative of a scalar. If they did, zero in one frame would imply zero in another. The lecture now slows down and proves that they do not.

\section{Testing the Ordinary Derivative of a Vector}

Let us follow the test carefully. Suppose we try to define
\begin{equation}
T_{mn}(x)\equiv \frac{\partial V_m(x)}{\partial x^n},
\end{equation}
and suppose, for the sake of argument, that this is the \(x\)-component of a genuine tensor. Then in \(y\)-coordinates it would have to transform as
\begin{align}
T_{mn}(y)
&=
\frac{\partial x^r}{\partial y^m}
\frac{\partial x^s}{\partial y^n}
T_{rs}(x) \\
&=
\frac{\partial x^r}{\partial y^m}
\frac{\partial x^s}{\partial y^n}
\frac{\partial V_r(x)}{\partial x^s}.
\end{align}
Using the chain rule on the last two factors, we obtain
\begin{equation}
T_{mn}(y)
=
\frac{\partial x^r}{\partial y^m}
\frac{\partial V_r(x)}{\partial y^n}.
\end{equation}

Now compare this with the actual derivative of the transformed vector components. Since
\begin{equation}
V_m(y)=\frac{\partial x^r}{\partial y^m}V_r(x),
\end{equation}
we find
\begin{align}
\frac{\partial V_m(y)}{\partial y^n}
&=
\frac{\partial}{\partial y^n}
\left(
\frac{\partial x^r}{\partial y^m}V_r(x)
\right) \\
&=
\frac{\partial x^r}{\partial y^m}\frac{\partial V_r(x)}{\partial y^n}
+
\frac{\partial}{\partial y^n}
\left(
\frac{\partial x^r}{\partial y^m}
\right)V_r(x).
\end{align}
The first term is exactly the transformed object we would have wanted. But there is also the second term,
\begin{equation}
\frac{\partial}{\partial y^n}
\left(
\frac{\partial x^r}{\partial y^m}
\right)V_r(x),
\end{equation}
and that term has no reason to vanish in curvilinear coordinates.

Therefore
\begin{equation}
\frac{\partial V_m(y)}{\partial y^n}\neq T_{mn}(y)
\end{equation}
in general. The ordinary derivative of a vector component does not transform as a tensor.

The lecture insists on lingering over this algebra, and rightly so. The important point is not just that the answer is no. It is that the product rule produces a new term when the derivative hits the coordinate-transformation factor. The failure is visible. One can see exactly where tensoriality is lost.

\subsection{Question \& Answer}

The natural question is: why does the ordinary derivative fail for a vector when it worked perfectly well for a scalar?

The answer is that a scalar does not carry along any changing basis information. A vector does. When we rewrite \(V_m(y)\) in terms of \(V_r(x)\), we must also rewrite the basis directions, and the factor \(\partial x^r/\partial y^m\) remembers that change of basis. Differentiating the vector therefore differentiates not only the components but also the coordinate frame itself. The extra term is the mathematical expression of the fact that the axes are moving around from point to point.

If the relation between the coordinates were linear, then \(\partial x^r/\partial y^m\) would be constant and its derivative would vanish. In that special case the trouble disappears. The failure is therefore not mysterious: it is the price we pay when the coordinates bend.

\section{Covariant Derivatives of Vectors and Tensors}

Once the extra term has appeared, the next move is almost forced. If ordinary differentiation fails by a very specific correction term, then perhaps the right derivative is ordinary differentiation plus that correction.

This is how the lecture introduces \(\Gamma\). One isolates the extra piece, gives it a name, and writes the corrected derivative as a new object. The board at this stage shows the covariant derivative of a lower-index vector in the form
\begin{figure}[h]
\centering
\includegraphics[width=0.72\textwidth]{lecture_05_figure_02.png}
\caption{Covariant derivative of a lower-index vector}
\end{figure}

\begin{equation}
\nabla_n V_m
=
\frac{\partial V_m}{\partial y^n}
+
\Gamma^r{}_{nm}V_r.
\end{equation}

\begin{remark}
At this point in the lecture the lower-index formula is written on the board with a plus sign in front of the connection term. We retain that local convention here in order to match the lecture's notation and the visible board work.
\end{remark}

The right-hand side is intended to be a genuine tensor, often denoted at this stage by \(T_{mn}\):
\begin{equation}
\frac{\partial V_m}{\partial y^n}
+
\Gamma^r{}_{nm}V_r
=
T_{mn}.
\end{equation}
This is the central repair. The word ``covariant'' in ``covariant derivative'' refers to the operator, not to the index type of the object being differentiated.

The lecture then broadens the rule. If a tensor has several lower indices, each lower index behaves like a vector index and contributes its own correction term. For a rank-two covariant tensor,
\begin{equation}
\nabla_p T_{mn}
=
\frac{\partial T_{mn}}{\partial y^p}
+
\Gamma^r{}_{pm}T_{rn}
+
\Gamma^r{}_{pn}T_{mr}.
\end{equation}
The classroom explanation is memorable: turn a blind eye to all indices except one, treat the visible index as if it belonged to a vector, write the corresponding correction, and then repeat the process index by index.

Later in the lecture the same rule is rewritten with the derivative index relabeled to \(p\), and the board returns to the one-index example:
\begin{figure}[h]
\centering
\includegraphics[width=0.72\textwidth]{lecture_05_figure_03.png}
\caption{Reindexed covariant-derivative formula}
\end{figure}

\begin{equation}
\nabla_p V_m
=
\frac{\partial V_m}{\partial y^p}
+
\Gamma^r{}_{pm}V_r.
\end{equation}

At precisely this point the audience raises the right bookkeeping question: what indices may be renamed, and what indices may not? The lecture's answer should be kept explicit.

\begin{itemize}
\item A repeated summation index is a dummy index and may be renamed.
\item A free index is not a dummy index and may not be interchanged with another free index or with a repeated one.
\item One may interchange indices only if one has independent symmetry information about the tensor in question.
\end{itemize}

The lecture later also writes the corresponding formula for an upper-index vector:
\begin{equation}
\nabla_m V^n
=
\frac{\partial V^n}{\partial y^m}
+
\Gamma^n{}_{mr}V^r.
\end{equation}
This form will be the useful one when we differentiate along a curve.

\section{A Concrete Example: Polar Coordinates}

After the index-heavy middle stretch, the lecture resets with an example that makes the geometry plain. In polar coordinates let
\begin{equation}
V_\theta=0,
\qquad
V_r=1.
\end{equation}
These components are constant. They are just \(0\) and \(1\) everywhere. But the vector field itself is not constant, because it points radially outward from the origin. At one point it points to the right, at another upward, at another to the left. Nobody looking at the arrows would call this a constant vector field.

So here is a decisive lesson: constant components do not imply a constant vector field. The covariant derivative is designed to detect the genuine geometric variation that the ordinary derivative of components can miss.

In the standard polar metric on the plane,
\begin{equation}
ds^2=dr^2+r^2\,d\theta^2,
\end{equation}
the nontrivial \(r\)-dependence sits in the \(\theta\theta\) part of the metric. That is exactly why Christoffel symbols appear even though the underlying space is flat. Flatness does not mean that the connection vanishes in every coordinate system. It means that there exist Cartesian coordinates in which the metric is \(\delta_{mn}\) and the Christoffel symbols vanish. In curvilinear coordinates on the same flat plane they do not vanish.

\section{Metric Compatibility and the Christoffel Symbols}

So far \(\Gamma\) has been introduced as whatever must be added to an ordinary derivative to restore tensoriality. The lecture now asks how \(\Gamma\) can be determined from the geometry itself.

The key observation is made first in flat Cartesian coordinates. There the metric is
\begin{equation}
g_{mn}=\delta_{mn},
\end{equation}
constant from point to point. If our covariant derivative is the right one, then the covariant derivative of the metric ought to vanish there. But the covariant derivative of the metric is supposed to be a tensor, so if it vanishes in one coordinate system, it must vanish in every coordinate system. This motivates the condition of metric compatibility:
\begin{equation}
\nabla_r g_{mn}=0.
\end{equation}

With the lecture's local convention for lower indices, this expands as
\begin{equation}
\frac{\partial g_{mn}}{\partial y^r}
+
\Gamma^s{}_{rm}g_{sn}
+
\Gamma^s{}_{rn}g_{ms}
=
0.
\end{equation}
This is now an equation for the connection in terms of the metric and its derivatives.

The lecture does not work through the full algebra on the board. The transcript in this stretch is also noisy. So here it is best to state the cleaned endpoint in its familiar normalized form:
\begin{equation}
\Gamma^A{}_{BC}
=
\frac{1}{2}g^{AD}
\left(
\partial_B g_{DC}
+
\partial_C g_{DB}
-
\partial_D g_{BC}
\right).
\end{equation}

This is the Christoffel symbol, also called the connection or affine connection. It is built from first derivatives of the metric. It is the object that makes covariant differentiation possible.

It is essential to say one more thing plainly: the Christoffel symbols are not tensors. The lecture gives the simplest reason. In flat space with Cartesian coordinates, the metric components are constant, so all Christoffel symbols vanish. But after a change to curvilinear coordinates in the same flat space, the metric acquires position dependence and the Christoffel symbols no longer vanish. A tensor that is zero in one frame would have to be zero in all frames. Therefore \(\Gamma^A{}_{BC}\) is not a tensor.

What the connection does have is a precise geometric role. At any point we can choose coordinates as close as possible to Cartesian coordinates. In such locally flattened coordinates the metric is \(\delta_{mn}\) at that point, and the first derivatives of the metric can also be made to vanish there. At that point the Christoffel symbols vanish, and the covariant derivative reduces to the ordinary derivative. This gives the right geometric interpretation: the covariant derivative tells us how a tensor varies when measured in coordinates that are, at the point of interest, as flat as possible.

\section{Differentiation Along Curves and Geodesics}

The last movement of the lecture shifts from pointwise differentiation to differentiation along a curve. Let the curve be parameterized by arc length \(s\). For a scalar field \(\phi\), the derivative along the curve is simply
\begin{equation}
\frac{d\phi}{ds}
=
\frac{\partial \phi}{\partial y^m}\frac{dy^m}{ds}.
\end{equation}
The factor
\begin{equation}
t^m\equiv \frac{dy^m}{ds}
\end{equation}
is the tangent vector to the curve. Because \(s\) is arc length, \(t^m\) is a unit tangent vector.

For a vector field the same idea must be corrected in the usual way: ordinary derivatives are replaced by covariant derivatives. Thus the covariant derivative of a vector along the curve is
\begin{equation}
\frac{D V^n}{Ds}
=
\nabla_m V^n\,\frac{dy^m}{ds}.
\end{equation}
Expanding \(\nabla_m V^n\) gives
\begin{equation}
\frac{D V^n}{Ds}
=
\frac{dV^n}{ds}
+
\Gamma^n{}_{mr}V^r\frac{dy^m}{ds}.
\end{equation}
The first term is the usual chain-rule derivative along the path; the second term compensates for the flopping around of the coordinates along that path.

Now comes the decisive step. Put \(V^n=t^n=dy^n/ds\). Then the covariant derivative of the tangent vector along the curve becomes
\begin{equation}
\frac{D t^n}{Ds}
=
\frac{d^2 y^n}{ds^2}
+
\Gamma^n{}_{mr}\frac{dy^m}{ds}\frac{dy^r}{ds}.
\end{equation}

In flat space with Cartesian coordinates this reduces to the ordinary derivative of the tangent vector. A straight line has constant tangent vector, so the result is zero. The lecture then asks the natural next question: can we still find curves in a genuinely curved space whose tangent vector is covariantly constant? The answer is yes.

\begin{definition}
A geodesic is a curve whose tangent vector is covariantly constant:
\begin{equation}
\frac{D t^n}{Ds}=0.
\end{equation}
Equivalently,
\begin{equation}
\frac{d^2 y^n}{ds^2}
+
\Gamma^n{}_{mr}\frac{dy^m}{ds}\frac{dy^r}{ds}
=
0.
\end{equation}
\end{definition}

This is the geodesic equation. It says that although a curved space has no globally straight lines in the Euclidean sense, it does have curves that are straightest possible. In coordinates that are locally as flat as possible, the tangent vector of a geodesic is constant for a little while, until the coordinates begin to bend again.

The lecture closes by telling us where this is going physically. In space-time, straight world-lines in special relativity are the motions of particles on which no force acts. In general relativity, free particles move on geodesics. The geodesic equation therefore plays the role of an equation of motion. The Christoffel symbols appear on the right-hand side like a force term, and since the Christoffel symbols are built from derivatives of the metric, the derivatives of the metric become the relativistic analog of gravitational force. Parallel transport and curvature are the next steps, but the essential apparatus is already in place.

\section{Summary}

The lecture develops one central idea from beginning to end: to write coordinate-independent laws, we must differentiate tensors in a way that produces tensors. Scalars cause no trouble, but vectors do, because their components remember the motion of the coordinate axes. The failure of the ordinary derivative is exposed explicitly by the product rule, and the repair is the covariant derivative, built using the connection.

From there the rest follows in a natural sequence. The metric compatibility condition determines the Christoffel symbols in terms of derivatives of the metric. The Christoffel symbols are not tensors, but they are exactly the objects needed to build tensorial derivatives. Contracting a covariant derivative with a tangent vector gives differentiation along a curve, and setting the covariant derivative of the tangent vector to zero defines geodesics. By the end of the lecture, the mathematical machinery is no longer an abstraction: it is already pointing toward the motion of particles in a gravitational field.